\section{Experimental Plan and Architecture}
\label{sec:experimental_plan}

\paragraph{Experimental Environment and Precision Standards}
All experiments execute without neural network gradient descent, backpropagation, or parameter optimization. The post-processing pipeline operates on frozen representations extracted in a single forward pass.
\begin{itemize}
\item Hardware Infrastructure: Experiments run on an Apple Silicon M3 Max workstation (36 GB unified memory, 30-core GPU, 14-core CPU) running macOS and PyTorch with Metal Performance Shaders (MPS) acceleration.
\item Numerical Precision: Representations, covariance accumulations, and transformations are strictly standardized on single-precision floating point (\texttt{torch.float32}).
\item Determinism: Random seeds across \texttt{torch}, \texttt{numpy}, and Python \texttt{random} are fixed to $2026$.
\end{itemize}

\paragraph{Evaluated Models and Architectural Properties}
Table~\ref{tab:model_architectures} summarizes the four evaluated sub-1B architectures.

\begin{table*}[t]
\centering
\small
\caption{Architectural specifications and operational properties of evaluated sub-1B models.}
\label{tab:model_architectures}
\resizebox{\textwidth}{!}{%
\begin{tabular}{ll cccc l}
\toprule
\textbf{Model Identifier} & \textbf{Architectural Tier} & \textbf{Parameters} & \textbf{Native $d$} & \textbf{Context Limit} & \textbf{MRL} & \textbf{Default Pooling Mode} \\
\midrule
\texttt{google/embeddinggemma-300m} & Dedicated Contrastive & 308M & 768 & 2,048 & Yes & Attention Mean Pooling \\
\texttt{Qwen/Qwen3-Embedding-0.6B} & Dedicated Contrastive & 600M & 1,024 & 32,768 & Yes & Final Non-Padding Token \\
\texttt{google/gemma-3-1b-pt} & Base Causal Language Model & 1.0B & 1,152 & 32,768 & No & Mean and Last Token Evaluated \\
\texttt{Qwen/Qwen3.5-0.8B-Base} & Base Causal Language Model & 0.8B & 1,024 & 262,144 & No & Mean and Last Token Evaluated \\
\bottomrule
\end{tabular}%
}
\end{table*}

\paragraph{Post-Processing Transformations}
Transformations are decoupled into full-dimension processing ($d \to d$) and dimensional compression ($d \to k$). Transformed representations are projected onto the unit hypersphere: $x' \leftarrow x' / \|x'\|_2$.
\begin{itemize}
\item \texttt{baseline}: Raw unit-normalized vectors: $x' = x / \|x\|_2$.
\item \texttt{standardization}: Coordinate-wise z-score scaling: $x'_j = (x_j - \mu_j) / \sigma_j$ \citep{timkey2021all}.
\item \texttt{r1}: Direct empirical mean subtraction: $x' = (x - \mu) / \|x - \mu\|_2$ \citep{ren2025correcting}.
\item \texttt{r2}: Mean-subspace projection deflation: $x' = (x - (x^\top \hat{\mu})\hat{\mu}) / \|x - (x^\top \hat{\mu})\hat{\mu}\|_2$, where $\hat{\mu} = \mu / \|\mu\|_2$.
\item \texttt{soft\_zca}: Regularized ZCA whitening: $W = U (\Lambda + \epsilon I)^{-1/2} U^\top$, $x' = (x - \mu)W / \|(x - \mu)W\|_2$, with spectral regularizer $\epsilon = 10^{-4}$ \citep{diera2024isotropy}.
\item \texttt{abtt\_D}: Principal component deflation removing $D \in \{1, 2, 3\}$ leading eigenvectors: $x' = (I - \sum_{j=1}^D u_j u_j^\top)(x - \mu)$ \citep{mu2018all}.
\item \texttt{rand} (Control): Deflating an arbitrary Gaussian unit vector $v \sim \mathcal{N}(0, I_d)$.
\item \texttt{mc} (Control): Mean centering without hyperspherical L2 renormalization: $x' = x - \mu$.
\item Compression ($d \to k, k \in \{512, 256, 128, 64\}$):
  \begin{itemize}
  \item \texttt{prefix}: Coordinate prefix slicing $x[:k] / \|x[:k]\|_2$.
  \item \texttt{pca}: Projection onto top-$k$ eigenvectors $(x - \mu)U_k / \|(x - \mu)U_k\|_2$ ($\gamma = 0$).
  \item \texttt{whitening}: Truncated Whitening-$k$ scaling components by $\lambda_j^{-1/2}$ ($\gamma = 1.0$) \citep{su2021whitening}.
  \item \texttt{spectemp}: Adaptive SNR-tempered spectral projection $W_k = U_k \text{diag}((\lambda_j + \epsilon)^{-\gamma(k)/2})$, where $\gamma(k) = \min(1.0, \text{SNR}(k) / \text{SNR}(k_{\text{knee}}))$ derived via the Kneedle algorithm ($S=0.5$) \citep{li2026spectral}.
  \item \texttt{random\_truncation} (Control): Uniform random subsampling of $k$ coordinates.
  \end{itemize}
\end{itemize}

\paragraph{Evaluation Metrics and Rationale}
\begin{itemize}
\item Retrieval: Evaluated using \texttt{pytrec-eval} against official qrels. Primary metric is \texttt{ndcg\_at\_10} (evaluating graded relevance order). Secondary metrics are \texttt{recall\_at\_100} (candidate coverage) and \texttt{mrr\_at\_10} (reciprocal rank of first relevant passage).
\item Similarity: Primary metric is Spearman rank correlation $\rho \times 100$ (measuring monotonic ranking agreement). Secondary metric is Pearson linear correlation $r \times 100$.
\item Intrinsic Geometry: Measured across corpus documents and pooled sentence representations: centroid norm ($\|\mu\|_2$), average pairwise cosine similarity across 2,000 sampled pairs, Maximum Explainable Variance of the leading eigenvalue ($\text{MEV} = \lambda_1 / \sum \lambda_j$), Normalized Intrinsic Dimensionality ($\text{NID} = H / \log d$), and IsoScore \citep{rudman2022isoscore}.
\end{itemize}

\paragraph{Viva Demonstration Scenario}
During the viva examination, the end-to-end post-processing pipeline is demonstrated interactively via CLI:
\begin{itemize}
\item Step 1: Execute representation extraction on an MTEB v2 benchmark split:
  \texttt{uv run scripts/encode.py --model google/gemma-3-1b-pt --task ArguAna --pooling mean\_pooling}.
\item Step 2: Apply the full transformation matrix in memory without model forward passes:
  \texttt{uv run scripts/transform.py --model google/gemma-3-1b-pt --task ArguAna --pooling mean\_pooling}.
\item Step 3: Run batched inner-product search and compute trec-eval ranking metrics in under two seconds:
  \texttt{uv run scripts/eval\_retrieval.py --model google/gemma-3-1b-pt --task ArguAna --pooling mean\_pooling}.
\item Step 4: Display live metric printout demonstrating the +18.4\% relative nDCG@10 gain from baseline (0.3050) to Soft-ZCA (0.3610).
\end{itemize}
